\documentclass[11pt]{article}
\usepackage[margin=0.85in]{geometry}
\usepackage[T1]{fontenc}
\usepackage{lmodern}
\usepackage{hyperref}
\title{Tactix: Current-Source H2 Proof-Contract Decision}
\author{Report to Andrew}
\date{25 September 2026}
\begin{document}
\maketitle
\vspace{-1.5em}

\section*{Outcome}
I replayed the exact 21-variable H2 public-API action program against the
latest Release solver objects. The binary native-DSR certificate is
byte-identical to the historical parent proof. With the canonical no-option
\texttt{dsr-trim} command, both backward and forward checking reject it. With
the same pinned binary's diagnostic \texttt{--delete-implied-units} option,
both directions verify and both decoded LSR elaborations pass independent
Trestle. No production route or solver rule was changed.

\section*{Evidence and interpretation}
The exact run is hwulab job 1722271 (two CPUs, 8 GiB; one second), archived
with its driver, hashes, proof and checker logs in the
\href{../../benchmarks/probes/am423-h2-current/}{AM423 probe}. The
\href{../../benchmarks/reports/2026-09-25-am423-current-h2-proof-contract.md}{technical report}
records each status and the proof hash. In no-option mode the checker
retains six requested deletions because their clauses are current global
unit-propagation reasons after an explicit root unit. Its later direct-SR
witness is therefore checked against a larger clause population than the
solver's. This explains the specific replay; it is not proof of a general
repair algorithm. A fixed deletion reorder or suppressing a legal public
action would not meet the project's generality and N1 requirements.

Seven current Release N1 fixtures passed, including representative public
rule instances, online composition, policy-independent Sweep/Reduce reach,
and Decide. They do not prove universal N1. The AM401, AM412, and AM414
Claude-owned counted campaigns remain active on immutable older pins; their
results cannot be transferred to this source. No performance claim is made.

\section*{Decision and next work}
The canonical proof contract needs a deliberate choice. Preserving the
no-option route requires a checker-visible proof-state/emission design for
implied-unit-reason deletions that also preserves every legal action. Permitting
\texttt{--delete-implied-units} requires changing the calculus/oracle contract,
replacing contract-specific negative controls, broad original-CNF proof and
model fuzz, and rerunning the retained gates; the option is not a monotone
relaxation. I requested that choice and have kept the no-option route pending
an answer. H2 and combined N1--N4 remain unverified. The requested full
5,000-second Tier-D run (four CPUs, 32 GiB, hwulab) remains downstream of
those gates; none was submitted.

\end{document}
